% Conclusion

We have presented an operator learning framework for delay differential equations
using Fourier Neural Operators.  Our contributions include: (i) a benchmark suite
of five DDE families spanning discrete and distributed delays, (ii) a robust data
generation pipeline with quality control and reproducibility guarantees, and
(iii) baseline experiments demonstrating the feasibility and current limitations
of FNO-based operator learning for DDEs.

The baseline FNO achieves 4.6--31.8$\times$ improvement over naive persistence across
all families, with median relative $L^2$ errors ranging from 0.077 (DistExp) to 0.612
(Linear2).  The difficulty ranking---distributed delays easiest, multi-delay systems
hardest---provides a useful guide for future architecture development.  Out-of-distribution
experiments reveal that generalization to unseen delay values and history function classes
remains a significant challenge, with gaps up to 7$\times$ and 11$\times$ respectively.

\paragraph{Limitations.}
Our baseline uses a single architecture (93k-parameter FNO) without per-family
hyperparameter tuning.  The benchmark currently covers five families; while these
span the major DDE categories, state-dependent delays and neutral-type DDEs are
not yet included.  The OOD splits test one axis of variation at a time rather than
joint shifts.

\paragraph{Future work.}
A scaling study across model sizes (small/medium/large), dataset sizes
(2k/10k/50k), and additional families (including neutral DDEs, chaotic
Mackey--Glass variants, and stiff Van der Pol systems) is currently underway.
Other promising directions include: comparison with alternative operator learning
architectures (DeepONets, Transformers), training with mixed history function
classes to improve OOD robustness, and extending the framework to
state-dependent and time-varying delays.

\paragraph{Three-phase benchmark extension.}
The follow-on benchmarking effort is structured in three phases.
\emph{Phase~A} (theorem-verification suite) runs three small controlled
targets---continuous-lag toy, expansive linear delay, contractive
tanh delay---against PlainMLP, MLP with cyclic-shift augmentation at
multiple budgets, LEMO, and an LEMO$_{\sigma}$ ablation, in order to
verify directly the continuous-lag equivariance theorem, the
augmentation covering-radius lower bound, and the LEMO$_{\sigma}$
rollout-stability corollary.  \emph{Phase~B} (core DDE benchmark)
evaluates the 11-model zoo listed in Section~\ref{sec:scope-frozen-vs-phase-b}
across seven families (DistExp, VdP, Hutchinson, Mackey--Glass,
Predator--Prey, DistUniform, Linear2) and four evaluation splits
(in-distribution, lag-shift OOD, horizon OOD, and either
continuous-$\tau$ OOD or history-amplitude OOD depending on family),
stratified by a lag-dependence strength score (LDS) to distinguish
lag-dominant from weakly lag-dependent regimes.  \emph{Phase~C}
(stability deep-dive) sweeps $\sigma \in \{0.5, 0.7, 0.9, 0.99, \text{unconstrained}\}$
on three representative families to characterize the
accuracy--stability frontier of the LEMO$_{\sigma}$ certified subclass.
The scaffolding for Phase~B is committed
(\texttt{configs/sweep\_phase\_b.yaml},
\texttt{slurm/sweep\_phase\_b.py}); the remaining blockers are
data generation for Mackey--Glass and Predator--Prey, and an
off-grid continuous-$\tau$ data pipeline.
